\cotefile{93}
\scanpage{179}
\unsure{vient|vent|revient} équivalemment d'exprimer le th.~3, la relation
\[
\boxed{\operatorname{res}\binom{\widetilde C\omega}{\widetilde f_1,\ldots,\widetilde f_n}
 =\operatorname{res}\binom{\omega}{f_1^p\ldots f_n^p}}
\]
Lorsque la dim relative est 1, cette formule s'écrit
\[
\begin{aligned}
\operatorname{res}\frac{\widetilde C\omega}{\widetilde f}
 &=\operatorname{res}\frac{\omega}{f^p}
 &\text{i.e.}\quad &\boxed{\operatorname{res}\widetilde C\unsure{\varpi|\overline\omega}=\operatorname{res}\varpi}\\
\overset{\rotatebox{90}{$=$}}{\inserted{\operatorname{res}\widetilde C\left(\frac{\omega}{f}\right)}}
 &&&\text{si on pose }\varpi=\omega/\unsure{\widetilde f|f|f^p}.
\end{aligned}
\]
et lorsque de plus $S$ est parfait i.e. $F_S:S\longrightarrow S$ un isom, de sorte que $C$ est défini, on peut l'écrire
\[
\begin{aligned}
\deleted{\left(\operatorname{res}\frac{C\omega}{f}\right)^p=\operatorname{res}\frac{\omega}{f^p}}
 &\qquad \boxed{\left(\operatorname{res} C\varpi\right)^p=\operatorname{res}\varpi}\\
\deleted{\overset{\rotatebox{90}{$=$}}{\operatorname{res} C\left(\frac{\omega}{f^p}\right)}}&
\end{aligned}
\]

\uline{Remarques} \textcircled{1} Il y a aussi un hom de Gysin $f_i:f_*(\Omega^{\deleted{\unsure{k}}i}_{X/S})\longrightarrow\Omega^i_{X^{p/S}/S}$, défini grâce à l'hom \inserted{de Gysin} pour $i=n$ et la formule de projection. Mais comme $f^\ast:\Omega^i_{X^{p/S}/S}\longrightarrow f_*\Omega^i_{X/S}$ est nul pour $i\ne0$, on conclut que $f_i=0$ si $i\ne n$. En particulier si $i\ne0$, l'hom $f_i$ ne coïncide \uline{pas} en général avec l'opération de Cartier, \inserted{sur les formes fermées}.

\textcircled{2} Il y a lieu de définir aussi l'opération de Cartier \deleted{\unsure{disons}} sur les modules dualisants relatif $\underline{\omega}_{X/S}$, \inserted{sans hyp. de lissité,} lorsque la formation de celui-ci commute au chgt. de base $S\xrightarrow{F_S}S$ \unsure{.} (p.ex. $S$ parfait i.e. $F_S$ un isom), \deleted{\unsure{et}} $X$ étant loc de type fini sur $S$ loc \unsure{noeth} disons. On trouve
\[
 C=f_{X/S\,*}:f_{X/S\,*}(\omega_{X/S})\longrightarrow
 \omega_{X^{p/S}/S}\simeq\widetilde\omega_{X/S}
\]
d'où, en appliquant $f_{X/S}^{-1\,*}$, \inserted{et composant avec $f_{X/S}'(\widetilde\omega_{X/S})\longrightarrow\widetilde\omega_{X/S}^{(p)}$,} un hom (pas linéaire, mais $\Lambda$-linéaire)
\[
 \omega_{X/S}\longrightarrow\omega_{X/S}^{(p)}
\]
% archive sheet number: 178
